\section{\texorpdfstring{Kategorie DG i kategorie $A_\infty$}{Kategorie DG i kategorie A-infty}}\label{cat:sec:dg-a-infty}

W kategorii homotopijnej kompleksów (rozdział~\ref{cat:sec:triangulated-derived}) zachowujemy jedynie klasy homotopii morfizmów kompleksów. Możemy zachować znacznie więcej informacji: wszystkie morfizmy jednorodne, ich różniczkę oraz składanie zgodne z regułą Leibniza. Tak otrzymujemy kategorię różniczkową z gradacją, w skrócie kategorię DG (ang.\ differential graded category). Jest to szczególny przypadek kategorii wzbogaconej (definicja~\ref{cat:def:enriched-category}). Kategorie $A_\infty$ pozwalają dodatkowo zastąpić ścisłą łączność przez spójny układ homotopii wyższych rzędów. W dalszym ciągu używamy gradacji kohomologicznej, czyli różniczek stopnia $+1$.

\subsection{Kompleksy i reguła znaków Koszula}

Niech $k$ będzie przemiennym pierścieniem z jedynką. Przez kompleks rozumiemy kompleks kohomologiczny $k$-modułów, a przez $|x|$ stopień elementu jednorodnego $x$.

\begin{definicja}[Iloczyn tensorowy kompleksów]\label{cat:def:dg-tensor-complex}
Dla kompleksów $V,W$ definiujemy
\[
 (V\otimes_k W)^n=\bigoplus_{p+q=n}V^p\otimes_k W^q,
 \qquad
 d(v\otimes w)=d_Vv\otimes w+(-1)^{|v|}v\otimes d_Ww.
\]
Jednostką jest kompleks $k$ skoncentrowany w stopniu zerowym. Symetria jest zadana wzorem
\[
 \tau_{V,W}(v\otimes w)=(-1)^{|v||w|}w\otimes v.
\]
\end{definicja}

\begin{stwierdzenie}[Symetryczna struktura monoidalna kompleksów]\label{cat:prop:dg-tensor-monoidal}
Powyższe konstrukcje wraz z kanonicznymi izomorfizmami łączności i jednostki wyznaczają symetryczną kategorię monoidalną $\mathrm{Ch}(k)$ kompleksów i morfizmów kompleksów.
\end{stwierdzenie}
\begin{proof}
Dla $|v|=p$ otrzymujemy
\[
\begin{aligned}
 d^2(v\otimes w)
 &=d_V^2v\otimes w+
 ((-1)^{p+1}+(-1)^p)d_Vv\otimes d_Ww
 +v\otimes d_W^2w=0.
\end{aligned}
\]
Różniczka jest dobrze określona na iloczynach tensorowych modułów, ponieważ oba składniki są $k$-dwuliniowe w ustalonych stopniach. Element sumy prostej ma skończony nośnik, zatem różniczka jest określona również dla kompleksów nieograniczonych.

Dla jednorodnych $v,w,z$ różniczki obu nawiasowań iloczynu tensorowego trzech kompleksów mają po kanonicznym utożsamieniu postać
\[
 d_Vv\otimes w\otimes z+
 (-1)^{|v|}v\otimes d_Ww\otimes z+
 (-1)^{|v|+|w|}v\otimes w\otimes d_Zz.
\]
Stąd izomorfizm łączności jest morfizmem kompleksów. Podobnie izomorfizmy jednostki zachowują różniczki. Dla $|w|=q$ mamy
\[
\begin{aligned}
 d\tau(v\otimes w)
 &=(-1)^{pq}d_Ww\otimes v+
   (-1)^{pq+q}w\otimes d_Vv,\\
 \tau d(v\otimes w)
 &=(-1)^{(p+1)q}w\otimes d_Vv+
   (-1)^{p+p(q+1)}d_Ww\otimes v.
\end{aligned}
\]
Współczynniki odpowiednich składników są równe, więc $\tau$ jest morfizmem kompleksów. Złożenie dwóch symetrii mnoży tensor przez $(-1)^{2pq}=1$. W obu diagramach sześciokątnych znaki są zgodne, ponieważ $p(q+r)=pq+pr$ i $(p+q)r=pr+qr$. Pięciokąt łączności oraz trójkąt jednostki są przemienne: na prostych tensorach wszystkie drogi jedynie zmieniają nawiasowanie lub usuwają czynnik jednostkowy. Proste tensory generują odpowiednie moduły, więc równości zachodzą wszędzie.
\end{proof}

Dla jednorodnych odwzorowań stopni $|f|,|g|$ między modułami z gradacją stosujemy regułę Koszula
\begin{equation}\label{cat:eq:dg-koszul-tensor}
 (f\otimes g)(v\otimes w)=(-1)^{|g||v|}f(v)\otimes g(w).
\end{equation}
Zastosowanie wzoru dwa razy daje
\[
 (f'\otimes g')(f\otimes g)
 =(-1)^{|g'||f|}(f'f\otimes g'g),
\]
ponieważ całkowity wykładnik po lewej wynosi
$|g||v|+|g'|(|v|+|f|)$, a po prawej
$|g'||f|+(|g|+|g'|)|v|$.

\subsection{Wewnętrzny kompleks morfizmów}

\begin{definicja}[Kompleks morfizmów]\label{cat:def:dg-hom-complex}
Dla kompleksów $V,W$ definiujemy
\[
 \underline{\Hom}_k(V,W)^n=
 \prod_{p\in\mathbb Z}\Hom_k(V^p,W^{p+n}).
\]
Dla morfizmu jednorodnego $f$ stopnia $n$ różniczka jest dana przez
\begin{equation}\label{cat:eq:dg-hom-differential}
 Df=d_W f-(-1)^n f d_V.
\end{equation}
\end{definicja}

W definicji używamy iloczynu, a nie sumy prostej: morfizm jednorodny może mieć niezerowe składowe w nieskończenie wielu stopniach. Nie zakładamy, że kompleksy są ograniczone.

\begin{stwierdzenie}[Różniczka morfizmów i składanie]\label{cat:prop:dg-hom-leibniz}
Kompleks morfizmów ma $D^2=0$. Dla jednorodnych $f:U\to V$ i $g:V\to W$ zachodzi
\[
 D(gf)=(Dg)f+(-1)^{|g|}g(Df).
\]
Identyczności są cyklami stopnia zerowego.
\end{stwierdzenie}
\begin{proof}
Dla $|f|=n$ obliczamy
\[
\begin{aligned}
 D^2f
 &=d_W(d_W f-(-1)^nf d_V)
    -(-1)^{n+1}(d_W f-(-1)^nf d_V)d_V\\
 &=d_W^2 f-
   \bigl((-1)^n+(-1)^{n+1}\bigr)d_W f d_V-fd_V^2=0.
\end{aligned}
\]
Jeżeli $|g|=m$, to
\[
\begin{aligned}
 (Dg)f+(-1)^m g(Df)
 &=d_Wgf-(-1)^m gd_Vf
   +(-1)^m gd_Vf-(-1)^{m+n}gfd_U\\
 &=D(gf).
\end{aligned}
\]
Wreszcie $D(\id_V)=d_V-d_V=0$.
\end{proof}

\begin{twierdzenie}[Sprzężenie tensor--Hom]\label{cat:thm:dg-tensor-hom}
Istnieje naturalny izomorfizm kompleksów
\[
 \underline{\Hom}_k(U\otimes_k V,W)
 \cong
 \underline{\Hom}_k(U,\underline{\Hom}_k(V,W)).
\]
W szczególności $\mathrm{Ch}(k)$ jest symetryczną kategorią monoidalnie domkniętą.
\end{twierdzenie}
\begin{proof}
Morfizmowi jednorodnemu $F$ stopnia $n$ przypisujemy $\Phi(F)$ przez
$\Phi(F)(u)(v)=F(u\otimes v)$. Dla $|u|=p$ odwzorowanie $\Phi(F)(u)$ ma stopień $n+p$. Odwrotnie, dla jednorodnego $G:U\to\underline{\Hom}_k(V,W)$ stopnia $n$ wzór $u\otimes v\mapsto G(u)(v)$ jest dwuliniowy w poszczególnych stopniach, więc określa morfizm $U\otimes_k V\to W$ stopnia $n$. Obie operacje są wzajemnie odwrotne na prostych tensorach, które generują moduły.

Dla jednorodnych $u,v$ mamy
\[
\begin{aligned}
 (D\Phi(F))(u)(v)
 &=d_W F(u\otimes v)
 -(-1)^{n+|u|}F(u\otimes d_Vv)
 -(-1)^nF(d_Uu\otimes v)\\
 &=\Phi(DF)(u)(v).
\end{aligned}
\]
Otrzymujemy izomorfizm kompleksów. Naturalność względem morfizmów kompleksów w każdej zmiennej wynika przez podstawienie do wzoru $\Phi(F)(u)(v)=F(u\otimes v)$. Biorąc cykle stopnia zerowego, otrzymujemy naturalną bijekcję zbiorów morfizmów kompleksów, która jest żądanym sprzężeniem.
\end{proof}

\subsection{Kategorie DG i funktory DG}

\begin{definicja}[Kategoria DG]\label{cat:def:dg-category}
Kategorią DG nad $k$ nazywamy kategorię wzbogaconą w $\mathrm{Ch}(k)$. Jawnie, składa się ona z klasy obiektów, kompleksów $\mathcal C(X,Y)$, morfizmów kompleksów
\[
 \mathcal C(Y,Z)\otimes_k\mathcal C(X,Y)
 \longrightarrow\mathcal C(X,Z),\quad g\otimes f\longmapsto gf
\]
oraz jednostek $\id_X\in Z^0\mathcal C(X,X)$, spełniających aksjomaty łączności i jednostki. Różniczka spełnia regułę Leibniza
\[
 d(gf)=(dg)f+(-1)^{|g|}g(df)
\]
dla jednorodnych morfizmów. Tu $Z^n K=\ker(d:K^n\to K^{n+1})$.
\end{definicja}

Reguła Leibniza jest dokładnie warunkiem, aby składanie było morfizmem kompleksów względem definicji~\ref{cat:def:dg-tensor-complex}.

\begin{przyklad}[Kategoria DG kompleksów]\label{cat:ex:dg-complexes}
Obiekty to kompleksy $k$-modułów, a kompleksy morfizmów są dane definicją~\ref{cat:def:dg-hom-complex}. Składanie jest zwykłym składaniem odwzorowań jednorodnych; jego stopnie dodają się. Stwierdzenie~\ref{cat:prop:dg-hom-leibniz} daje zgodność składania z różniczką i domkniętość jednostek. Łączność i prawa jednostki wynikają z odpowiednich własności składania funkcji. Otrzymaną kategorię oznaczamy przez $\mathcal{C}_{\mathrm{dg}}(k)$.
\end{przyklad}

\begin{przyklad}[Algebra DG jako kategoria jednoobiektowa]\label{cat:ex:dg-one-object}
Jednoobiektowa kategoria DG jest tym samym co unitalna łączna algebra DG $A$: jej jedyny kompleks morfizmów to $A$, składanie to mnożenie, a jednostka to $1\in Z^0(A)$. W przeciwnym kierunku każda taka algebra określa kategorię DG o jednym obiekcie. Aksjomaty obu struktur są identyczne po tych utożsamieniach. Na przykład algebra z gradacją i różniczką zerową jest algebrą DG, jeśli jej mnożenie zachowuje stopnie.
\end{przyklad}

\begin{definicja}[Funktor DG]\label{cat:def:dg-functor}
Funktor DG $F:\mathcal C\to\mathcal D$ składa się z odwzorowania obiektów i morfizmów kompleksów
\[
 F_{X,Y}:\mathcal C(X,Y)\longrightarrow\mathcal D(FX,FY)
\]
zachowujących składanie i jednostki. Złożenie funktorów DG jest określone przez składanie odpowiednich odwzorowań obiektów i kompleksów morfizmów.
\end{definicja}

\subsection{Kategorie cykli i kohomologii}

\begin{definicja}[Kategorie \texorpdfstring{$Z^0(\mathcal C)$ i $H^0(\mathcal C)$}{cykli i kohomologii}]\label{cat:def:dg-h-zero}
Dla kategorii DG $\mathcal C$ definiujemy kategorie o tych samych obiektach, co $\mathcal C$, i zbiorach morfizmów
\[
 Z^0(\mathcal C)(X,Y)=Z^0\mathcal C(X,Y),\qquad
 H^0(\mathcal C)(X,Y)=H^0\mathcal C(X,Y).
\]
Składanie w pierwszej kategorii jest odziedziczone, a w drugiej dane przez $[g][f]=[gf]$.
\end{definicja}

\begin{stwierdzenie}[Dobre określenie składania]\label{cat:prop:dg-cohomology-composition}
Konstrukcje $Z^0$ i $H^0$ są funktorialne względem funktorów DG. Ogólniej grupy $H^n\mathcal C(X,Y)$ tworzą kategorię z gradacją, ze składaniem stopnia zerowego.
\end{stwierdzenie}
\begin{proof}
Dla cykli jednorodnych $g,f$ reguła Leibniza daje $d(gf)=0$. Jeżeli $f=da$ i $g$ jest cyklem, to $d(ga)=(-1)^{|g|}gda$, więc $gf$ jest brzegiem. Jeżeli $g=db$ i $f$ jest cyklem, to $d(bf)=(db)f$, więc również $gf$ jest brzegiem. Z dwuliniowości wynika niezależność klasy $[gf]$ od obu reprezentantów. Łączność i jednostki przechodzą do klas, a stopień złożenia wynosi $|g|+|f|$. Funktor DG zachowuje cykle i brzegi, gdyż komutuje z różniczką, oraz zachowuje składanie i jednostki, więc indukuje żądane funktory.
\end{proof}

\begin{stwierdzenie}[Odzyskanie kategorii homotopijnej]\label{cat:prop:dg-homotopy-category}
Zachodzą kanoniczne izomorfizmy kategorii
\[
 Z^0(\mathcal C_{\mathrm{dg}}(k))\cong\mathrm{Ch}(k),
 \qquad
 H^0(\mathcal C_{\mathrm{dg}}(k))\cong K(k),
\]
gdzie $K(k)$ oznacza kategorię homotopijną kompleksów $k$-modułów z rozdziału~\ref{cat:sec:triangulated-derived}.
\end{stwierdzenie}
\begin{proof}
Morfizm stopnia zerowego jest cyklem dokładnie wtedy, gdy $d_Wf=fd_V$, czyli jest morfizmem kompleksów. Dla $h$ stopnia $-1$ mamy $Dh=d_Wh+hd_V$. Różnica $f-g$ dwóch morfizmów kompleksów jest brzegiem dokładnie wtedy, gdy $f-g=d_Wh+hd_V$, co jest definicją homotopii (równanie~\eqref{cat:eq:cochain-homotopy}). Oba utożsamienia są identycznością na obiektach, są bijekcjami na zbiorach morfizmów i zachowują składanie przez stwierdzenie~\ref{cat:prop:dg-cohomology-composition}.
\end{proof}

\begin{stwierdzenie}[Wyższe kohomologie kompleksu morfizmów]\label{cat:prop:dg-hom-shift}
Dla wszystkich $n\in\mathbb Z$ istnieje naturalny izomorfizm $k$-modułów
\[
 H^n\underline{\Hom}_k(V,W)
 \cong \Hom_{K(k)}(V,W[n]),
\]
gdzie $W[n]^p=W^{p+n}$ i $d_{W[n]}=(-1)^n d_W$.
\end{stwierdzenie}
\begin{proof}
Przyjmujemy konwencję przesunięcia z definicji~\ref{cat:def:complex-shift}. Jednorodne $f$ stopnia $n$ można traktować jako morfizm stopnia zerowego $V\to W[n]$. Warunek $Df=0$ oznacza $d_Wf=(-1)^nf d_V$, równoważnie $d_{W[n]}f=fd_V$. Jeżeli $a$ ma stopień $n-1$, to
\[
 Da=d_Wa+(-1)^n a d_V
 =(-1)^n(d_{W[n]}a+a d_V).
\]
Zatem brzegi stopnia $n$ odpowiadają morfizmom homotopijnym z zerem, ponieważ mnożenie homotopii przez $(-1)^n$ jest odwracalne. Otrzymujemy izomorfizm na ilorazach. Naturalność względem morfizmów kompleksów stopnia zerowego wynika z tego, że utożsamienie nie zmienia składowych odwzorowania i komutuje ze składaniem. Morfizmy homotopijne indukują te same odwzorowania na kohomologii przez regułę Leibniza.
\end{proof}

\begin{przyklad}[Kompleks ściągalny i jego endomorfizmy]\label{cat:ex:dg-contractible}
Niech $V$ będzie kompleksem ściągalnym, czyli istnieje $h:V\to V$ stopnia $-1$ z $Dh=\id_V$. Dla każdego kompleksu $W$ oba kompleksy $\underline{\Hom}_k(W,V)$ i $\underline{\Hom}_k(V,W)$ są ściągalne. W pierwszym przypadku definiujemy homotopię $T(f)=hf$. Reguła Leibniza daje $D(hf)=f-hDf$, a zatem $DT+TD=\id$. W drugim przypadku dla $|g|=n$ kładziemy $T(g)=(-1)^ngh$. Wtedy
\[
 D(T(g))=(-1)^n(Dg)h+g,\qquad
 T(Dg)=(-1)^{n+1}(Dg)h,
\]
więc znowu $DT+TD=\id$. W szczególności wszystkie ich kohomologie znikają. To mocniejsze stwierdzenie niż sam fakt, że $V$ jest obiektem zerowym w kategorii homotopijnej.
\end{przyklad}

\subsection{Kategoria przeciwna DG i iloczyn tensorowy kategorii DG}

\begin{definicja}[Kategoria przeciwna DG]\label{cat:def:dg-opposite}
Kategoria $\mathcal C^{\mathrm{op}}$ ma te same obiekty co $\mathcal C$, kompleksy $\mathcal C^{\mathrm{op}}(X,Y)=\mathcal C(Y,X)$ z niezmienioną różniczką, a dla $f:X\to Y$, $g:Y\to Z$ rozumianych jako morfizmy przeciwne definiujemy
\[
 g\circ_{\mathrm{op}} f=(-1)^{|f||g|}fg,
\]
gdzie po prawej stronie składamy w $\mathcal C$.
\end{definicja}

\begin{stwierdzenie}[Poprawność kategorii przeciwnej]\label{cat:prop:dg-opposite}
Definicja~\ref{cat:def:dg-opposite} określa kategorię DG. Zachodzą $(\mathcal C^{\mathrm{op}})^{\mathrm{op}}=\mathcal C$ i $H^0(\mathcal C^{\mathrm{op}})=H^0(\mathcal C)^{\mathrm{op}}$ przy kanonicznych utożsamieniach.
\end{stwierdzenie}
\begin{proof}
Niech $|f|=p$, $|g|=q$, $|h|=r$. Oba nawiasowania przeciwnego złożenia trzech morfizmów dają $fgh$ ze znakiem $(-1)^{pq+pr+qr}$, gdyż $(p+q)r+pq=p(q+r)+qr$. Jednostki mają stopień zerowy, więc nie wprowadzają znaku. Reguła Leibniza przyjmuje postać
\[
\begin{aligned}
 d(g\circ_{\mathrm{op}}f)
 &=(-1)^{pq}\bigl((df)g+(-1)^pf(dg)\bigr),\\
 (dg)\circ_{\mathrm{op}}f+(-1)^qg\circ_{\mathrm{op}}(df)
 &=(-1)^{p(q+1)}f(dg)
 +(-1)^{q+(p+1)q}(df)g.
\end{aligned}
\]
Wykładniki odpowiadających sobie składników są równe modulo $2$. Przejście dwa razy do kategorii przeciwnej mnoży każdy znak przez samego siebie. W stopniu zerowym znak wynosi $1$, co daje ostatnie utożsamienie.
\end{proof}

\begin{definicja}[Iloczyn tensorowy kategorii DG]\label{cat:def:dg-category-tensor}
Dla kategorii DG $\mathcal C,\mathcal D$ nad tym samym $k$ obiektami $\mathcal C\otimes_k\mathcal D$ są pary $(X,Y)$. Kompleksy morfizmów definiujemy przez
\[
 (\mathcal C\otimes_k\mathcal D)((X,Y),(X',Y'))
 =\mathcal C(X,X')\otimes_k\mathcal D(Y,Y').
\]
Składanie na jednorodnych prostych tensorach to
\[
 (f'\otimes g')(f\otimes g)
 =(-1)^{|g'||f|}(f'f\otimes g'g).
\]
Jednostką obiektu $(X,Y)$ jest $\id_X\otimes\id_Y$.
\end{definicja}

\begin{stwierdzenie}[Poprawność iloczynu tensorowego]\label{cat:prop:dg-category-tensor}
Powyższa konstrukcja określa kategorię DG.
\end{stwierdzenie}
\begin{proof}
Składanie jest złożeniem morfizmów kompleksów: najpierw używamy symetrii z definicji~\ref{cat:def:dg-tensor-complex}, aby przestawić środkowe czynniki w tensorze $f'\otimes g'\otimes f\otimes g$, a następnie składamy w obu kategoriach. Zatem jest morfizmem kompleksów. Dla trzech kolejnych tensorów $f\otimes g,f'\otimes g',f''\otimes g''$ oba nawiasowania składania dają $(f''f'f)\otimes(g''g'g)$ ze znakiem
\[
 (-1)^{|g'||f|+|g''||f|+|g''||f'|}.
\]
Wynika to z dodawania stopni i dwukrotnego zastosowania wzoru z definicji. Jednostki są cyklami, bo oba ich czynniki są cyklami; ich stopień jest zerowy, więc składanie z jednostką nie zmienia znaku. Powyższe rachunki rozszerzają się liniowo, ponieważ proste tensory generują odpowiednie moduły.
\end{proof}

\subsection{Quasi-równoważności}

\begin{definicja}[Quasi-równoważność kategorii DG]\label{cat:def:dg-quasi-equivalence}
Funktor DG $F:\mathcal C\to\mathcal D$ nazywamy quasi-równoważnością, jeżeli:
\begin{enumerate}
 \item $F_{X,Y}$ jest quasi-izomorfizmem kompleksów dla każdej pary $X,Y$, czyli indukuje izomorfizmy we wszystkich stopniach kohomologii;
 \item $H^0(F):H^0(\mathcal C)\to H^0(\mathcal D)$ jest istotnie surjektywny na obiektach.
\end{enumerate}
\end{definicja}

\begin{stwierdzenie}[Quasi-równoważność a kategoria homotopijna]\label{cat:prop:dg-quasi-equivalence-h0}
Quasi-równoważność indukuje równoważność kategorii $H^0$. Złożenie quasi-równoważności jest quasi-równoważnością. Odwrotna implikacja, od równoważności $H^0$ do quasi-równoważności, jest na ogół fałszywa.
\end{stwierdzenie}
\begin{proof}
Pierwszy warunek definicji w stopniu zerowym oznacza pełność i wierność $H^0(F)$; drugi oznacza istotną surjektywność. Są to kryteria równoważności kategorii. Złożenie quasi-izomorfizmów jest quasi-izomorfizmem, gdyż we wszystkich stopniach indukuje złożenie izomorfizmów. Złożenie funktorów istotnie surjektywnych jest istotnie surjektywne: obiekt w ostatniej kategorii jest izomorficzny z obrazem obiektu pośredniego, który z kolei jest izomorficzny z obrazem obiektu początkowego.

Dla kontrprzykładu niech $k$ będzie niezerowym pierścieniem przemiennym, a $A=k\oplus k\varepsilon$ z $|\varepsilon|=1$, $\varepsilon^2=0$ i różniczką zerową. Inkluzja $k\to A$ określa funktor między jednoobiektowymi kategoriami DG. W stopniu zerowym jest izomorfizmem, więc indukuje izomorfizm kategorii $H^0$. W stopniu pierwszym indukuje jednak $0\to k\varepsilon$, co nie jest izomorfizmem. Funktor nie jest zatem quasi-równoważnością.
\end{proof}

